\documentclass[11pt]{article}

\usepackage[margin=1in]{geometry}
\usepackage{amsmath, amssymb, amsfonts, bm}
\usepackage{booktabs}
\usepackage{enumitem}
\usepackage{graphicx}
\usepackage{hyperref}
\usepackage{xcolor}
\graphicspath{{../../artifacts/}}

\title{Posterior Update Math for Toy and Identity Benchmarks}
\author{kinematic-classifier-sandbox}
\date{\today}

\begin{document}

\maketitle

\section{Purpose}

This document is the beginning of a formal mathematical write-up for the two
active benchmark families in the repository:
\begin{itemize}[leftmargin=1.5em]
    \item \texttt{toy\_1d.py}: a class-matched latent-state filter bank
    \item \texttt{identity\_1d.py}: a direct speed-identity classifier over
    \texttt{bike}, \texttt{horse}, and \texttt{car}
\end{itemize}

The aim is to connect the implemented code paths to explicit posterior update
equations and to show how those updates produce the downstream confusion
matrices, feature summaries, and entropy traces. The generic recursive update
language is standard Bayesian filtering notation, while the repo-specific value
is the common evidence contract and artifact traceability built around it
\cite{bishop2006,kalman1960}.

This document also serves a second purpose. It is a methodology note for how
the repository is being used as a classification study system. That means the
write-up needs to cover not only equations, but also the exploratory structure
around those equations:
\begin{itemize}[leftmargin=1.5em]
    \item how corpora are generated and stressed,
    \item how features and filters are interpreted as evidence providers,
    \item how posterior histories are turned into diagnostic artifacts,
    \item and what the current 1D benchmarks do and do not yet prove.
\end{itemize}

Primary implementation surfaces covered explicitly by this note:
\begin{itemize}[leftmargin=1.5em]
    \item \texttt{toy\_1d.py}
    \item \texttt{identity\_1d.py}
    \item \texttt{posterior\_explainer.py}
    \item \texttt{identity\_posterior\_explainer.py}
    \item \texttt{bayesian\_walkthroughs.py}
    \item posterior-oriented artifact writers in \texttt{artifacts.py}
\end{itemize}

\section{Methodology Framing}

The mathematical story in this repository is not just ``a classifier exists.''
The more important claim is methodological:
\begin{itemize}[leftmargin=1.5em]
    \item a corpus layer produces observations, truth, class labels, and
    scenario labels
    \item a feature or filter layer turns those observations into evidence
    \item a posterior layer accumulates that evidence over time
    \item an analysis layer measures confusion, separability, prior
    sensitivity, and entropy
    \item a reporting layer turns those standard artifacts into tables and
    plots
\end{itemize}

So the 1D benchmarks are being used as a proving ground for a broader
classification-and-filtering methodology. The key structural question is not
only whether the current 1D benchmarks work, but whether the interfaces
between:
\begin{equation}
    \text{corpus} \rightarrow \text{features or filters} \rightarrow
    \text{evidence} \rightarrow \text{posterior} \rightarrow \text{metrics}
\end{equation}
are generic enough that later 3D work mostly means changing the state
representation, the feature extractors, and the dynamics or measurement models,
not rewriting the whole experiment and analysis machinery.

\subsection{Layered Methodology View}

The repository is converging toward a layered interpretation:
\begin{align}
    \text{corpus layer}
    &\rightarrow \text{feature or filter layer}
    \rightarrow \text{evidence layer} \\
    &\rightarrow \text{posterior layer}
    \rightarrow \text{analysis layer}
    \rightarrow \text{reporting layer}. \nonumber
\end{align}

Those layers play different roles:
\begin{itemize}[leftmargin=1.5em]
    \item The corpus layer defines what is observed and what truth is
    available.
    \item The feature or filter layer transforms raw trajectories into
    evidence-bearing summaries, innovations, or residuals.
    \item The posterior layer accumulates evidence over time.
    \item The analysis layer measures confusion, identifiability, prior
    sensitivity, adequacy, and calibration.
    \item The reporting layer renders the results as plots, tables, and
    reproducible artifact bundles.
\end{itemize}

The 1D studies are valuable because they let those layers be stress-tested in
a controlled setting before the geometry becomes more complicated in 3D.

\subsection{Exploratory Versus Operational Roles}

The current 1D algorithms should be read as exploratory methodology tools,
not as final operational trackers. They are designed to answer questions such
as:
\begin{itemize}[leftmargin=1.5em]
    \item Which classes are separable from the available evidence?
    \item Which confusions are feature problems versus posterior-update
    problems?
    \item How much of a decision is driven by likelihood versus prior?
    \item When does a model-based filter add value over static features?
\end{itemize}

That distinction matters because a benchmark can be useful even when it is not
yet highly accurate. A weak benchmark may still reveal:
\begin{itemize}[leftmargin=1.5em]
    \item an inadequate feature family,
    \item a biased corpus,
    \item an over-separable or under-separable class pair,
    \item or a posterior mechanism that is too brittle.
\end{itemize}

\section{What Is Random and What Is Being Inferred}

The cleanest reading of the repository is:
\begin{itemize}[leftmargin=1.5em]
    \item $s_i$ denotes class $i$
    \item $z_t$ denotes the measurement at time $t$
    \item $x_t$ denotes a latent kinematic state only in the toy benchmark
    \item $p(s_i \mid z_{1:t})$ is the posterior class probability after seeing
    measurements up to time $t$
\end{itemize}

The two benchmark families do not infer the same latent object:
\begin{itemize}[leftmargin=1.5em]
    \item \texttt{toy\_1d.py} maintains a latent state posterior inside each
    class-conditioned filter
    \item \texttt{identity\_1d.py} does not maintain a latent Kalman state and
    instead updates class probabilities directly from observed speed
\end{itemize}

\section{Generic Recursive Posterior Update}

For either benchmark, let $s_i$ denote class $i$ and let $z_{1:t}$ denote the
measurement history up to time $t$. The recursive class update is
\begin{equation}
    p(s_i \mid z_{1:t}) \propto p(s_i \mid z_{1:t-1}) \, L_i(z_t, z_{1:t-1}),
\end{equation}
with normalization
\begin{equation}
    p(s_i \mid z_{1:t}) =
    \frac{p(s_i \mid z_{1:t-1}) \, L_i(z_t, z_{1:t-1})}
    {\sum_j p(s_j \mid z_{1:t-1}) \, L_j(z_t, z_{1:t-1})}.
\end{equation}

In log space, define
\begin{equation}
    \ell_{i,t} = \log p(s_i \mid z_{1:t-1}) + \log L_i(z_t, z_{1:t-1}),
\end{equation}
so that
\begin{equation}
    p(s_i \mid z_{1:t}) =
    \frac{\exp(\ell_{i,t})}{\sum_j \exp(\ell_{j,t})}.
\end{equation}

\subsection{Generic Interpretation of the Update}

The clean abstraction behind the repository is:
\begin{equation}
    \underbrace{\text{prior over classes}}_{\text{posterior at } t-1}
    +
    \underbrace{\text{new evidence}}_{\log L_i}
    \Longrightarrow
    \underbrace{\text{posterior at } t}_{\text{normalized score}}.
\end{equation}

This means very different method families can still be compared through the
same posterior machinery:
\begin{itemize}[leftmargin=1.5em]
    \item pointwise methods produce direct class evidence from a single
    measurement or feature vector
    \item windowed methods produce class evidence from short histories
    \item sequential Bayesian accumulators produce cumulative evidence
    \item Kalman banks produce evidence from innovations and state residuals
\end{itemize}

The posterior updater should not need to know where the evidence came from,
only that it is compatible log-likelihood information.

\subsection{Evidence-Provider Interpretation}

One meta-level point of the current work is that the method families can be
viewed through a shared evidence-provider abstraction. In schematic form,
\begin{equation}
    \mathcal{E}(\text{observation history}, \text{features}, \text{model state})
    \rightarrow
    \{\log L_i\}_{i=1}^K,
\end{equation}
where $\mathcal{E}$ may be:
\begin{itemize}[leftmargin=1.5em]
    \item a pointwise measurement scorer,
    \item a windowed feature scorer,
    \item a cumulative Bayesian accumulator,
    \item or a Kalman-style innovation and residual scorer.
\end{itemize}

This is one of the central methodological claims of the repo: different
classifiers and filters can be compared as long as they emit compatible class
evidence histories and posterior histories.

\subsection{Implementation Mapping}

The mathematical objects in this note correspond to explicit code surfaces in
the current repository.

\begin{center}
\begin{tabular}{p{0.28\textwidth} p{0.28\textwidth} p{0.34\textwidth}}
\toprule
\textbf{Concept} & \textbf{Primary implementation surface} & \textbf{Role} \\
\midrule
Gaussian interval mass & \texttt{toy\_1d.gaussian\_interval\_probability} & Soft symmetric envelope probability for speed and acceleration \\
Innovation log likelihood & \texttt{toy\_1d.\_innovation\_log\_likelihood} & Textbook Gaussian measurement-fit term \\
Toy recursive class bank & \texttt{toy\_1d.run\_class\_bank} & Per-class prediction, composite scoring, posterior update, and feature probabilities \\
Identity recursive benchmark & \texttt{identity\_1d.run\_identity\_benchmark} & Direct speed evidence, soft validity, short-window mode logic, and posterior update \\
Toy posterior artifact generation & \texttt{posterior\_explainer.py} & Step-level success, failure, comparison, and margin-trace diagnostics \\
Identity posterior artifact generation & \texttt{identity\_posterior\_explainer.py} & Boundary-failure and posterior-term diagnostics for the identity benchmark \\
\bottomrule
\end{tabular}
\end{center}

So the document should be read as implementation-linked mathematics rather than
as a free-standing theoretical note.

\subsection{PDF Terms Versus Probability-Mass Terms}

A key implementation detail is that the repository uses two different types of
probabilistic terms:
\begin{itemize}[leftmargin=1.5em]
    \item density terms, which score how well an exact observed value fits a
    class model
    \item region-probability terms, which integrate Gaussian mass over a valid
    region of the state or measurement space
\end{itemize}

So the answer to ``are we chopping the Gaussian probability space?'' is:
\begin{itemize}[leftmargin=1.5em]
    \item mostly no for the main innovation or speed-shape terms
    \item yes, softly, for the envelope or validity terms
\end{itemize}

\begin{figure}[htbp]
    \centering
    \includegraphics[width=\textwidth]{probability_primitives.png}
    \caption{Simple explanatory charts for the probability primitives used in
    the 1D classifiers: innovation density, symmetric soft limits, one-sided
    soft limits, and prior-to-posterior updating.}
\end{figure}

\begin{figure}[htbp]
    \centering
    \includegraphics[width=\textwidth]{posterior_update_math.png}
    \caption{Rendered overview of the posterior update structure for the toy and
    identity 1D classifiers.}
\end{figure}

\section{Toy 1D Benchmark}

\subsection{State and Measurement Model}

The toy benchmark uses a latent state
\begin{equation}
    x_t =
    \begin{bmatrix}
        p_t \\
        v_t \\
        a_t
    \end{bmatrix},
\end{equation}
with scalar position measurement
\begin{equation}
    z_t = H x_t + \nu_t,
    \qquad
    H = \begin{bmatrix} 1 & 0 & 0 \end{bmatrix},
    \qquad
    \nu_t \sim \mathcal{N}(0, R).
\end{equation}

Each class maintains its own predicted state mean and covariance,
$\hat{x}_{i,t}^{-}$ and $P_{i,t}^{-}$, and its own updated moments after
measurement incorporation.

\subsection{How the Toy Classifier Is Set Up}

The toy benchmark is a class bank. For each class, the code:
\begin{enumerate}[leftmargin=1.8em]
    \item maintains a separate latent state mean and covariance
    \item predicts that state forward with class-specific dynamics
    \item scores the new position measurement under that class
    \item adds soft plausibility terms on velocity, acceleration, direction,
    oscillation, and short-window behavior
    \item updates the posterior weight of that class
\end{enumerate}

So this is not a single shared tracker followed by an external classifier. It
is a joint recursive filter-bank style classifier.

\subsection{Toy 1D Class Semantics}

Inside the current 1D framework, the toy class bank is meant to separate
different motion behaviors rather than static identities:
\begin{itemize}[leftmargin=1.5em]
    \item \texttt{brake}: deceleration and stopping pressure
    \item \texttt{coast}: low-drive, low-structure motion
    \item \texttt{drift}: sustained reverse or backtracking behavior
    \item \texttt{maneuver}: oscillatory or regime-switching behavior
    \item \texttt{powered}: persistent positive drive and higher-speed motion
    \item \texttt{unknown}: motions that are intentionally hard to anchor to a
    cleaner named class
\end{itemize}

The important point is that these are not only end labels. They correspond to
different expected evidence patterns over time, which is why the benchmark also
tracks phases, transient summaries, terminal summaries, and feature-vs-class
diagnostics.

\subsection{How Well the Toy 1D Algorithms Work Inside This Framework}

The toy benchmark should be evaluated as a methodology probe, not just a final
scoreboard. In the current artifact set:
\begin{itemize}[leftmargin=1.5em]
    \item overall trajectory accuracy is about $0.63$
    \item transient accuracy is about $0.46$
    \item terminal accuracy is about $0.54$
    \item \texttt{drift}, \texttt{powered}, and \texttt{unknown} are strong
    classes
    \item \texttt{brake} and especially \texttt{coast} remain weak
    \item transient and terminal summaries are materially different, which
    confirms that a whole-track label is hiding regime changes
\end{itemize}

This is useful because the benchmark is doing more than producing one final
confusion matrix. It exposes:
\begin{itemize}[leftmargin=1.5em]
    \item class confusion across full trajectories,
    \item transient and terminal confusion separately,
    \item phase confusion,
    \item feature-versus-class confusion,
    \item and posterior entropy collapse over time.
\end{itemize}

The entropy traces are especially useful. The mean class entropy falls from
roughly $1.48$ at the first step to about $0.19$ by step $20$. That does not
mean every decision is correct, but it does mean the posterior machinery is
accumulating decisive evidence rather than remaining uniformly ambiguous.

The toy benchmark is therefore proving that the framework can distinguish:
\begin{itemize}[leftmargin=1.5em]
    \item feature failures,
    \item likelihood-shape failures,
    \item prior or mode-mixture failures,
    \item and class-definition or corpus-design failures.
\end{itemize}

\subsection{Innovation Likelihood}

For class $s_i$, the innovation is
\begin{equation}
    r_{i,t} = z_t - H \hat{x}_{i,t}^{-},
    \qquad
    S_{i,t} = H P_{i,t}^{-} H^\top + R.
\end{equation}

The innovation log likelihood is
\begin{equation}
    \log L^{\text{dyn}}_{i,t}
    =
    -\frac{1}{2}
    \left[
        \log(2\pi S_{i,t}) + \frac{r_{i,t}^2}{S_{i,t}}
    \right].
\end{equation}

This is the most textbook piece of the toy classifier. It is a Gaussian
density term on the innovation, not an interval probability.

\subsection{Soft Envelope Terms}

The benchmark uses covariance-aware envelope probabilities rather than hard
thresholds. For example,
\begin{equation}
    \log L^{\text{speed}}_{i,t}
    =
    \lambda_v \log\!\Big(
        \epsilon + P(|v_t| \leq v_{\max,i} \mid s_i)
    \Big),
\end{equation}
and
\begin{equation}
    \log L^{\text{accel}}_{i,t}
    =
    \lambda_a \log\!\Big(
        \epsilon + P(|a_t| \leq a_{\max,i} \mid s_i)
    \Big).
\end{equation}

\subsection{How the Toy Model Cuts Probability Space}

The toy model softly cuts probability space through Gaussian CDF-based region
probabilities. For example, instead of using only a density like
\begin{equation}
    \mathcal{N}(v_t; \mu_v, \sigma_v^2),
\end{equation}
it also asks for the posterior mass inside a class-valid interval:
\begin{equation}
    P(-v_{\max,i} \leq v_t \leq v_{\max,i} \mid s_i).
\end{equation}

This is a soft geometric cut:
\begin{itemize}[leftmargin=1.5em]
    \item if most of the velocity posterior lies inside the valid interval, the
    class is rewarded
    \item if much of that posterior lies outside, the class is penalized
\end{itemize}

The same logic applies to acceleration.

\subsection{Composite Toy Score}

The implemented toy score can be written schematically as
\begin{align}
    \log L_{i,t}
    &=
    \log L^{\text{dyn}}_{i,t}
    + \log L^{\text{speed}}_{i,t}
    + \log L^{\text{accel}}_{i,t} \\
    &\quad
    + \log L^{\text{behavior}}_{i,t}
    + \log L^{\text{obs}}_{i,t}
    + \log L^{\text{mode}}_{i,t} \nonumber \\
    &\quad
    - \mathbf{1}\{s_i = \text{unknown}\}\gamma_{\text{unknown}}. \nonumber
\end{align}

Here:
\begin{itemize}[leftmargin=1.5em]
    \item $\log L^{\text{behavior}}_{i,t}$ scores latent velocity,
    acceleration, direction, and oscillation consistency.
    \item $\log L^{\text{obs}}_{i,t}$ scores finite-difference observed
    kinematics from recent measurements.
    \item $\log L^{\text{mode}}_{i,t}$ is a within-class mode mixture used
    especially for \texttt{brake} and \texttt{maneuver}.
\end{itemize}

\subsection{Toy Update Loop in Implementation Terms}

Inside the implementation, \texttt{run\_class\_bank} performs the toy update
loop in the following order:
\begin{enumerate}[leftmargin=1.8em]
    \item predict latent state moments for each class,
    \item evaluate the innovation likelihood,
    \item evaluate soft speed and acceleration region terms,
    \item evaluate behavior, observed-kinematics, and within-class mode terms,
    \item add the previous log class prior,
    \item normalize across classes with a log-sum-exp step,
    \item record posterior weights, feature probabilities, phase labels, and
    entropy.
\end{enumerate}

That final recording step matters. The benchmark is designed so that the same
recursive pass which produces a class posterior also produces:
\begin{itemize}[leftmargin=1.5em]
    \item per-step likelihood-term breakdowns,
    \item per-step posterior entropy,
    \item per-step feature probabilities,
    \item and per-step phase diagnostics.
\end{itemize}

This is why the downstream artifact layer can render not only final confusion
matrices but also success and failure walkthroughs with term-level explanations.

\subsection{What the Toy Posterior Represents}

The toy classifier is carrying two layers of inference:
\begin{itemize}[leftmargin=1.5em]
    \item inside each class: a state posterior $p(x_t \mid s_i, z_{1:t})$
    \item across classes: a class posterior $p(s_i \mid z_{1:t})$
\end{itemize}

That distinction matters because the benchmark can simultaneously hold:
\begin{itemize}[leftmargin=1.5em]
    \item a posterior class weight for \texttt{drift}
    \item a latent mean and covariance for the \texttt{drift} state estimate
\end{itemize}

\section{Identity Benchmark}

\subsection{Measurement Model}

The identity benchmark uses a direct scalar speed observation
\begin{equation}
    z_t = \text{observed speed in mph}.
\end{equation}

There is no latent Kalman state in this benchmark. The posterior update is
driven by a class-conditioned speed-shape model and several short-window
history terms.

\subsection{How the Identity Classifier Is Set Up}

At each step, the identity classifier takes one observed speed and asks:
\begin{itemize}[leftmargin=1.5em]
    \item how well does this speed fit the \texttt{bike} speed model?
    \item how well does it fit the \texttt{horse} speed model?
    \item how well does it fit the \texttt{car} speed model?
\end{itemize}

It then adjusts those instantaneous fits with:
\begin{itemize}[leftmargin=1.5em]
    \item a soft upper-speed validity term
    \item a running-history term
    \item a short-window mode term
    \item a recent-dynamics term
\end{itemize}

So the identity model is a recursive class-evidence updater, not a latent-state
tracker.

\subsection{Identity 1D Class Semantics}

The identity benchmark is deliberately simpler. Its class semantics are closer
to static speed-regime identities:
\begin{itemize}[leftmargin=1.5em]
    \item \texttt{bike}: lower-speed envelope with occasional surging behavior
    \item \texttt{horse}: mid-speed envelope, especially near the upper horse
    limit
    \item \texttt{car}: higher-speed envelope with persistent push or overspeed
    relative to horse
\end{itemize}

That simpler setup is methodologically useful because it isolates what direct
observation-space evidence can do before latent-state filtering enters the
picture.

\subsection{How Well the Identity Benchmark Works Inside This Framework}

The identity benchmark plays a different methodological role from the toy
benchmark. It is simpler dynamically, but it is valuable because it isolates
class evidence from direct speed observations and short-window history terms.

In the current artifact set:
\begin{itemize}[leftmargin=1.5em]
    \item overall accuracy is about $0.88$
    \item transient accuracy is about $0.86$
    \item terminal accuracy is about $0.89$
    \item \texttt{car} and \texttt{horse} are strong
    \item the main boundary pressure remains near the
    \texttt{bike}--\texttt{horse} overlap
    \item posterior entropy collapses quickly, which is expected for a simpler
    direct-speed identity problem
\end{itemize}

Methodologically, this benchmark helps answer a different question from toy:
\begin{itemize}[leftmargin=1.5em]
    \item how much separation is available from direct observation-space
    evidence alone?
    \item when is a class boundary mostly an envelope problem?
    \item when is it a short-history shape problem?
\end{itemize}

That is why the identity path is useful as a contrast case. It shows how far a
recursive evidence model can go without a full latent-state filter bank.

\subsection{Composite Identity Score}

For class $s_i$, the current code is well approximated by
\begin{equation}
    \log L_{i,t}
    =
    \log L^{\text{speed-shape}}_{i,t}
    + \log L^{\text{valid}}_{i,t}
    + \log L^{\text{history}}_{i,t}
    + \log L^{\text{mode}}_{i,t}
    + \log L^{\text{dyn-shape}}_{i,t}.
\end{equation}

The main pieces are:
\begin{align}
    \log L^{\text{speed-shape}}_{i,t}
    &= \log \mathcal{N}\!\left(z_t; \mu_i, \sigma_i^2 + \sigma_{\text{obs}}^2\right), \\
    \log L^{\text{valid}}_{i,t}
    &= 1.4 \log P\!\left(z_t \leq v_{\max,i} + \delta_i\right), \\
    \log L^{\text{history}}_{i,t}
    &= 0.45 \log \mathcal{N}\!\left(\bar{z}_{1:t}; \mu_i, \sigma_i^2 + \sigma_{\text{hist},t}^2\right).
\end{align}

The mode and dynamics terms capture recent-window structure such as cruise
regimes, border modes, surge behavior, or persistent push.

\subsection{Identity Update Loop in Implementation Terms}

Inside \texttt{run\_identity\_benchmark}, the identity benchmark performs a
lighter-weight recursive pass:
\begin{enumerate}[leftmargin=1.8em]
    \item read the current observed speed,
    \item evaluate the direct speed-shape log density for each class,
    \item evaluate the one-sided validity mass term,
    \item evaluate running-history, mode-shape, and recent-dynamics terms,
    \item add the previous class prior in log space,
    \item normalize across classes,
    \item record posterior, detected features, and entropy.
\end{enumerate}

So the identity benchmark shares the same posterior structure as toy while
removing the latent-state machinery. That makes it a useful control case for
the broader methodology.

\subsection{How the Identity Model Cuts Probability Space}

The identity model does not cut probability space through the main
\texttt{speed-shape} term. That term is a Gaussian density on the observed
speed. The soft cut appears in the validity term:
\begin{equation}
    P(z_t \leq v_{\max,i} + \delta_i).
\end{equation}

So instead of using a hard rule
\[
    z_t \leq v_{\max,i},
\]
the classifier uses a one-sided Gaussian probability mass term. This means that
a noisy observation slightly above the nominal class limit does not force the
class posterior to zero; it is only penalized.

\subsection{What the Identity Posterior Represents}

The identity benchmark carries only the class posterior
\begin{equation}
    p(s_i \mid z_{1:t}),
\end{equation}
with no separate latent-state posterior.

Compared with toy:
\begin{itemize}[leftmargin=1.5em]
    \item toy has class posterior plus per-class state posterior
    \item identity has class posterior only
\end{itemize}

\section{Worked Numeric Walkthrough Procedure}

For a given step $t$ and class $s_i$, the most direct numeric walkthrough is:
\begin{align}
    \ell_{i,t}^{\text{prior}} &= \log p(s_i \mid z_{1:t-1}), \\
    \ell_{i,t}^{\text{total}} &= \log L_i(z_t, z_{1:t-1}), \\
    \ell_{i,t}^{\text{num}} &= \ell_{i,t}^{\text{prior}} + \ell_{i,t}^{\text{total}}, \\
    \log Z_t &= \log \sum_j \exp(\ell_{j,t}^{\text{num}}), \\
    p(s_i \mid z_{1:t}) &= \exp\!\left(\ell_{i,t}^{\text{num}} - \log Z_t\right).
\end{align}

This sequence corresponds exactly to the rendered artifact
\texttt{posterior\_numeric\_walkthrough.(md, svg, png)} in the repo.

\begin{figure}[htbp]
    \centering
    \includegraphics[width=\textwidth]{posterior_numeric_walkthrough.png}
    \caption{Numeric substitution of the posterior equations for one toy failure
    case and one identity boundary failure case.}
\end{figure}

\section{Features, Confusion, and Entropy}

\subsection{Features as Evidence-Carrying Objects}

In this repository, features are not the final prediction by themselves. Their
role is to carry evidence. A useful feature therefore needs to be understood
through at least four questions:
\begin{itemize}[leftmargin=1.5em]
    \item What geometric or temporal behavior does it summarize?
    \item Is it instantaneous, windowed, cumulative, or model-residual based?
    \item How sensitive is it to duration, sample count, and noise?
    \item Does it transfer to higher-dimensional settings directly, through a
    norm, or only through a frame-specific projection?
\end{itemize}

This is why the repository now emphasizes feature confusion matrices and
feature-taxonomy work. The important methodological point is:
\begin{equation}
    \text{feature} \longrightarrow \text{evidence term}
    \longrightarrow \text{posterior behavior}
    \longrightarrow \text{metric and confusion artifact}.
\end{equation}

If a class is weak, the failure might arise from:
\begin{itemize}[leftmargin=1.5em]
    \item poor feature excitation,
    \item poor feature selectivity,
    \item a bad evidence model on otherwise useful features,
    \item or a corpus that is too easy or too hard for the declared class pair.
\end{itemize}

The feature probabilities are diagnostic quantities rather than the class
posterior itself. They are used to build feature-vs-class matrices and feature
confusion counts.

\subsection{Feature Families and Meta-Level Interpretation}

Within the current 1D studies, features can be grouped by role rather than only
by name:
\begin{itemize}[leftmargin=1.5em]
    \item instantaneous features summarize one step or a local state estimate
    \item windowed or shape features summarize short temporal structure
    \item cumulative-style features summarize what has happened over a larger
    prefix of the track
    \item model-residual features summarize mismatch between observation and a
    dynamics model
    \item derivative features summarize change rates and regime switches
\end{itemize}

That grouping matters because the feature studies in the repo are not only
asking ``which feature is strongest?'' They are asking:
\begin{itemize}[leftmargin=1.5em]
    \item which family carries the most separability,
    \item which family is brittle to noise or duration,
    \item which family overlaps too strongly across classes,
    \item and which family is most likely to transfer to 3D through norms,
    projections, residuals, or frame-aware geometry.
\end{itemize}

For example, many scalar 1D feature ideas should not be copied to 3D verbatim,
but lifted through structurally similar objects:
\begin{itemize}[leftmargin=1.5em]
    \item scalar speed features become speed-magnitude or component-projection
    features,
    \item scalar acceleration features become acceleration norms or
    frame-projected accelerations,
    \item and scalar innovation features become Mahalanobis innovation norms or
    direction-aware residual summaries.
\end{itemize}

\subsection{From Feature Evidence to Study Artifacts}

The repository is increasingly treating feature work as a study pipeline rather
than a one-off classifier detail:
\begin{equation}
    \text{feature extractor}
    \rightarrow
    \text{feature table}
    \rightarrow
    \text{evidence behavior}
    \rightarrow
    \text{confusion, overlap, PCA, and adequacy artifacts}.
\end{equation}

That is why feature information now appears in several distinct artifact
families:
\begin{itemize}[leftmargin=1.5em]
    \item feature-vs-class confusion heatmaps,
    \item feature precision, recall, and lift summaries,
    \item feature-set inspection bundles,
    \item PCA summaries and loading plots,
    \item and corpus adequacy or excitation reports.
\end{itemize}

In other words, a feature is not only something the classifier consumes. It is
also something the study framework audits and compares.

Class entropy at step $t$ is
\begin{equation}
    H_t = - \sum_i p(s_i \mid z_{1:t}) \log p(s_i \mid z_{1:t}).
\end{equation}

This quantity is useful for distinguishing:
\begin{itemize}[leftmargin=1.5em]
    \item confident correct classification,
    \item confident incorrect classification,
    \item persistent ambiguity across classes.
\end{itemize}

\subsection{What the Current 1D Feature Studies Prove}

The current 1D feature studies already support several meta-level claims:
\begin{itemize}[leftmargin=1.5em]
    \item feature sets can be compared independently of any one classifier
    family,
    \item pairwise class difficulty can be separated from posterior-update
    mechanics,
    \item and the same feature-space summaries support adequacy audits, PCA,
    prior sensitivity, and recommendation tooling.
\end{itemize}

This matters because it is the beginning of a generic methodology layer. In a
future 3D setting, the specific scalar features will change, but the study
questions should remain the same:
\begin{itemize}[leftmargin=1.5em]
    \item Which feature families excite the intended classes?
    \item Which class pairs remain intrinsically hard?
    \item Which confusions are evidence-limited versus model-limited?
\end{itemize}

\section{Methodological Conclusions From the Current 1D Work}

The current 1D benchmark family supports the following conclusions:
\begin{enumerate}[leftmargin=1.8em]
    \item The posterior machinery is already rich enough to compare
    pointwise-style evidence, windowed evidence, and model-based evidence under
    a common Bayesian update story.
    \item The feature and confusion artifacts are strong enough to diagnose why
    a class fails, not just that it fails.
    \item The toy filter-bank benchmark and the identity direct-evidence
    benchmark are complementary:
    \begin{itemize}[leftmargin=1.5em]
        \item toy stresses latent-state and mode-mixture reasoning,
        \item identity stresses direct observation-space evidence and boundary
        calibration.
    \end{itemize}
    \item The current numbers are strong enough to support exploratory claims,
    but not yet strong enough to claim finished operational 1D solutions:
    \begin{itemize}[leftmargin=1.5em]
        \item toy still has weak \texttt{brake} and \texttt{coast} behavior,
        \item identity still has informative \texttt{bike}--\texttt{horse}
        and \texttt{horse}--\texttt{car} boundary pressure.
    \end{itemize}
    \item The repo is already more than a single toy classifier, but it is
    still using 1D studies as a methodology proving ground rather than claiming
    a finished dimension-agnostic framework.
\end{enumerate}

\subsection{What ``Working Well'' Means in This Framework}

Inside this repository, a method is not judged only by top-line accuracy. It is
working well when it supports the broader study loop:
\begin{itemize}[leftmargin=1.5em]
    \item the evidence terms are interpretable,
    \item posterior histories can be inspected step by step,
    \item entropy and confidence behave coherently,
    \item class and feature confusions point to actionable model or corpus
    changes,
    \item and the outputs fit the common analysis and artifact machinery.
\end{itemize}

By that standard, the current 1D systems are already productive methodology
instruments. They are not yet the end-state algorithms, but they are
successfully exercising the classification, evidence, posterior, and reporting
framework that later 3D work is expected to reuse.

\subsection{Implementation Consequence for Future Documentation}

Because the repository now has a clearer separation between feature studies,
posterior studies, filtering studies, and showcase packaging, future
documentation should continue to connect three levels explicitly:
\begin{itemize}[leftmargin=1.5em]
    \item the mathematical object,
    \item the implementation surface that computes it,
    \item and the artifact family that exposes its behavior.
\end{itemize}

That is the simplest way to keep the methodology legible as the project moves
from 1D exploratory work toward more generic and later higher-dimensional
studies.

\begin{figure}[htbp]
    \centering
    \includegraphics[width=0.48\textwidth]{toy_1d_feature_confusion.png}
    \hfill
    \includegraphics[width=0.48\textwidth]{identity_1d_feature_confusion.png}
    \caption{Rendered feature-to-class confusion heatmaps for the toy benchmark
    (left) and the identity benchmark (right).}
\end{figure}

\section{Next Extensions}

This starter document is intentionally narrow. Good next additions would be:
\begin{itemize}[leftmargin=1.5em]
    \item explicit notation for every toy behavior subterm,
    \item explicit notation for the identity mode-mixture terms,
    \item a short feature-taxonomy appendix aligned to the repository registry,
    \item a dimensional-lift appendix describing how these 1D feature ideas
    transfer to 3D through norms, projections, and residuals,
    \item one worked numerical example copied directly from a benchmark run,
    \item and a short evidence-provider section connecting pointwise,
    windowed, accumulator, and Kalman methods under one posterior contract.
\end{itemize}

\section{References}
\begin{thebibliography}{9}
\bibitem{bishop2006}
C. M. Bishop.
\newblock \emph{Pattern Recognition and Machine Learning}.
\newblock Springer, 2006.

\bibitem{kalman1960}
R. E. Kalman.
\newblock A new approach to linear filtering and prediction problems.
\newblock \emph{Transactions of the ASME}, 82D:35--45, 1960.
\end{thebibliography}

\end{document}
